\IfFileExists{stacks-project.cls}{%
\documentclass{stacks-project}
}{%
\documentclass{amsart}
}

% For dealing with references we use the comment environment
\usepackage{verbatim}
\newenvironment{reference}{\comment}{\endcomment}
%\newenvironment{reference}{}{}
\newenvironment{slogan}{\comment}{\endcomment}
\newenvironment{history}{\comment}{\endcomment}

% For commutative diagrams we use Xy-pic
\usepackage[all]{xy}

% We use 2cell for 2-commutative diagrams.
\xyoption{2cell}
\UseAllTwocells

% We use multicol for the list of chapters between chapters
\usepackage{multicol}

% This is generally recommended for better output
\usepackage{lmodern}
\usepackage[T1]{fontenc}

% For cross-file-references
\usepackage{xr-hyper}

% Package for hypertext links:
\usepackage{hyperref}

% For any local file, say "hello.tex" you want to link to please
% use \externaldocument[hello-]{hello}
\externaldocument[introduction-]{introduction}
\externaldocument[conventions-]{conventions}
\externaldocument[sets-]{sets}
\externaldocument[categories-]{categories}
\externaldocument[topology-]{topology}
\externaldocument[sheaves-]{sheaves}
\externaldocument[sites-]{sites}
\externaldocument[stacks-]{stacks}
\externaldocument[fields-]{fields}
\externaldocument[algebra-]{algebra}
\externaldocument[brauer-]{brauer}
\externaldocument[homology-]{homology}
\externaldocument[derived-]{derived}
\externaldocument[simplicial-]{simplicial}
\externaldocument[more-algebra-]{more-algebra}
\externaldocument[smoothing-]{smoothing}
\externaldocument[modules-]{modules}
\externaldocument[sites-modules-]{sites-modules}
\externaldocument[injectives-]{injectives}
\externaldocument[cohomology-]{cohomology}
\externaldocument[sites-cohomology-]{sites-cohomology}
\externaldocument[dga-]{dga}
\externaldocument[dpa-]{dpa}
\externaldocument[sdga-]{sdga}
\externaldocument[hypercovering-]{hypercovering}
\externaldocument[schemes-]{schemes}
\externaldocument[constructions-]{constructions}
\externaldocument[properties-]{properties}
\externaldocument[morphisms-]{morphisms}
\externaldocument[coherent-]{coherent}
\externaldocument[divisors-]{divisors}
\externaldocument[limits-]{limits}
\externaldocument[varieties-]{varieties}
\externaldocument[topologies-]{topologies}
\externaldocument[descent-]{descent}
\externaldocument[perfect-]{perfect}
\externaldocument[more-morphisms-]{more-morphisms}
\externaldocument[flat-]{flat}
\externaldocument[groupoids-]{groupoids}
\externaldocument[more-groupoids-]{more-groupoids}
\externaldocument[etale-]{etale}
\externaldocument[chow-]{chow}
\externaldocument[intersection-]{intersection}
\externaldocument[pic-]{pic}
\externaldocument[weil-]{weil}
\externaldocument[adequate-]{adequate}
\externaldocument[dualizing-]{dualizing}
\externaldocument[duality-]{duality}
\externaldocument[discriminant-]{discriminant}
\externaldocument[derham-]{derham}
\externaldocument[local-cohomology-]{local-cohomology}
\externaldocument[algebraization-]{algebraization}
\externaldocument[curves-]{curves}
\externaldocument[resolve-]{resolve}
\externaldocument[models-]{models}
\externaldocument[functors-]{functors}
\externaldocument[equiv-]{equiv}
\externaldocument[pione-]{pione}
\externaldocument[etale-cohomology-]{etale-cohomology}
\externaldocument[proetale-]{proetale}
\externaldocument[relative-cycles-]{relative-cycles}
\externaldocument[more-etale-]{more-etale}
\externaldocument[trace-]{trace}
\externaldocument[crystalline-]{crystalline}
\externaldocument[spaces-]{spaces}
\externaldocument[spaces-properties-]{spaces-properties}
\externaldocument[spaces-morphisms-]{spaces-morphisms}
\externaldocument[decent-spaces-]{decent-spaces}
\externaldocument[spaces-cohomology-]{spaces-cohomology}
\externaldocument[spaces-limits-]{spaces-limits}
\externaldocument[spaces-divisors-]{spaces-divisors}
\externaldocument[spaces-over-fields-]{spaces-over-fields}
\externaldocument[spaces-topologies-]{spaces-topologies}
\externaldocument[spaces-descent-]{spaces-descent}
\externaldocument[spaces-perfect-]{spaces-perfect}
\externaldocument[spaces-more-morphisms-]{spaces-more-morphisms}
\externaldocument[spaces-flat-]{spaces-flat}
\externaldocument[spaces-groupoids-]{spaces-groupoids}
\externaldocument[spaces-more-groupoids-]{spaces-more-groupoids}
\externaldocument[bootstrap-]{bootstrap}
\externaldocument[spaces-pushouts-]{spaces-pushouts}
\externaldocument[spaces-chow-]{spaces-chow}
\externaldocument[groupoids-quotients-]{groupoids-quotients}
\externaldocument[spaces-more-cohomology-]{spaces-more-cohomology}
\externaldocument[spaces-simplicial-]{spaces-simplicial}
\externaldocument[spaces-duality-]{spaces-duality}
\externaldocument[formal-spaces-]{formal-spaces}
\externaldocument[restricted-]{restricted}
\externaldocument[spaces-resolve-]{spaces-resolve}
\externaldocument[formal-defos-]{formal-defos}
\externaldocument[defos-]{defos}
\externaldocument[cotangent-]{cotangent}
\externaldocument[examples-defos-]{examples-defos}
\externaldocument[algebraic-]{algebraic}
\externaldocument[examples-stacks-]{examples-stacks}
\externaldocument[stacks-sheaves-]{stacks-sheaves}
\externaldocument[criteria-]{criteria}
\externaldocument[artin-]{artin}
\externaldocument[quot-]{quot}
\externaldocument[stacks-properties-]{stacks-properties}
\externaldocument[stacks-morphisms-]{stacks-morphisms}
\externaldocument[stacks-limits-]{stacks-limits}
\externaldocument[stacks-cohomology-]{stacks-cohomology}
\externaldocument[stacks-perfect-]{stacks-perfect}
\externaldocument[stacks-introduction-]{stacks-introduction}
\externaldocument[stacks-more-morphisms-]{stacks-more-morphisms}
\externaldocument[stacks-geometry-]{stacks-geometry}
\externaldocument[moduli-]{moduli}
\externaldocument[moduli-curves-]{moduli-curves}
\externaldocument[examples-]{examples}
\externaldocument[exercises-]{exercises}
\externaldocument[guide-]{guide}
\externaldocument[desirables-]{desirables}
\externaldocument[coding-]{coding}
\externaldocument[obsolete-]{obsolete}
\externaldocument[fdl-]{fdl}
\externaldocument[index-]{index}

% Theorem environments.
%
\theoremstyle{plain}
\newtheorem{theorem}[subsection]{Theorem}
\newtheorem{proposition}[subsection]{Proposition}
\newtheorem{lemma}[subsection]{Lemma}

\theoremstyle{definition}
\newtheorem{definition}[subsection]{Definition}
\newtheorem{example}[subsection]{Example}
\newtheorem{exercise}[subsection]{Exercise}
\newtheorem{situation}[subsection]{Situation}

\theoremstyle{remark}
\newtheorem{remark}[subsection]{Remark}
\newtheorem{remarks}[subsection]{Remarks}

\numberwithin{equation}{subsection}

% Macros
%
\def\lim{\mathop{\mathrm{lim}}\nolimits}
\def\colim{\mathop{\mathrm{colim}}\nolimits}
\def\Spec{\mathop{\mathrm{Spec}}}
\def\Hom{\mathop{\mathrm{Hom}}\nolimits}
\def\Ext{\mathop{\mathrm{Ext}}\nolimits}
\def\SheafHom{\mathop{\mathcal{H}\!\mathit{om}}\nolimits}
\def\SheafExt{\mathop{\mathcal{E}\!\mathit{xt}}\nolimits}
\def\Sch{\mathit{Sch}}
\def\Mor{\mathop{\mathrm{Mor}}\nolimits}
\def\Ob{\mathop{\mathrm{Ob}}\nolimits}
\def\Sh{\mathop{\mathit{Sh}}\nolimits}
\def\NL{\mathop{N\!L}\nolimits}
\def\CH{\mathop{\mathrm{CH}}\nolimits}
\def\proetale{{pro\text{-}\acute{e}tale}}
\def\etale{{\acute{e}tale}}
\def\QCoh{\mathit{QCoh}}
\def\Ker{\mathop{\mathrm{Ker}}}
\def\Im{\mathop{\mathrm{Im}}}
\def\Coker{\mathop{\mathrm{Coker}}}
\def\Coim{\mathop{\mathrm{Coim}}}

% Boxtimes
%
\DeclareMathSymbol{\boxtimes}{\mathbin}{AMSa}{"02}

%
% Macros for moduli stacks/spaces
%
\def\QCohstack{\mathcal{QC}\!\mathit{oh}}
\def\Cohstack{\mathcal{C}\!\mathit{oh}}
\def\Spacesstack{\mathcal{S}\!\mathit{paces}}
\def\Quotfunctor{\mathrm{Quot}}
\def\Hilbfunctor{\mathrm{Hilb}}
\def\Curvesstack{\mathcal{C}\!\mathit{urves}}
\def\Polarizedstack{\mathcal{P}\!\mathit{olarized}}
\def\Complexesstack{\mathcal{C}\!\mathit{omplexes}}
% \Pic is the operator that assigns to X its picard group, usage \Pic(X)
% \Picardstack_{X/B} denotes the Picard stack of X over B
% \Picardfunctor_{X/B} denotes the Picard functor of X over B
\def\Pic{\mathop{\mathrm{Pic}}\nolimits}
\def\Picardstack{\mathcal{P}\!\mathit{ic}}
\def\Picardfunctor{\mathrm{Pic}}
\def\Deformationcategory{\mathcal{D}\!\mathit{ef}}
\usepackage{cleveref}
\crefname{originallemma}{original lemma}{original lemmas}
\Crefname{originallemma}{Original lemma}{Original lemmas}
\makeatletter
\newlabel{descent-lemma-standard-covering-Cech}{{03FI}{source}{}{sourcepointer-1}{}}
\newlabel{descent-lemma-standard-covering-Cech@cref}{{[originallemma][0][]03FI}{[1][0][]source}{}{}{}}
\makeatother
\crefname{originalproposition}{original proposition}{original propositions}
\Crefname{originalproposition}{Original proposition}{Original propositions}
\makeatletter
\newlabel{descent-proposition-same-cohomology-quasi-coherent}{{03DW}{source}{}{sourcepointer-2}{}}
\newlabel{descent-proposition-same-cohomology-quasi-coherent@cref}{{[originalproposition][0][]03DW}{[1][0][]source}{}{}{}}
\makeatother
\crefname{originallemma}{original lemma}{original lemmas}
\Crefname{originallemma}{Original lemma}{Original lemmas}
\makeatletter
\newlabel{lemma-adequate-flat}{{06UW}{source}{}{sourcepointer-3}{}}
\newlabel{lemma-adequate-flat@cref}{{[originallemma][0][]06UW}{[1][0][]source}{}{}{}}
\makeatother
\crefname{originallemma}{original lemma}{original lemmas}
\Crefname{originallemma}{Original lemma}{Original lemmas}
\makeatletter
\newlabel{lemma-adequate-fpqc}{{06VL}{source}{}{sourcepointer-4}{}}
\newlabel{lemma-adequate-fpqc@cref}{{[originallemma][0][]06VL}{[1][0][]source}{}{}{}}
\makeatother
\crefname{originallemma}{original lemma}{original lemmas}
\Crefname{originallemma}{Original lemma}{Original lemmas}
\makeatletter
\newlabel{lemma-adequate-local}{{06VI}{source}{}{sourcepointer-5}{}}
\newlabel{lemma-adequate-local@cref}{{[originallemma][0][]06VI}{[1][0][]source}{}{}{}}
\makeatother
\crefname{originallemma}{original lemma}{original lemmas}
\Crefname{originallemma}{Original lemma}{Original lemmas}
\makeatletter
\newlabel{sites-cohomology-lemma-cech-vanish-collection}{{03F9}{source}{}{sourcepointer-6}{}}
\newlabel{sites-cohomology-lemma-cech-vanish-collection@cref}{{[originallemma][0][]03F9}{[1][0][]source}{}{}{}}
\makeatother
\title{An adequate module is quasi-coherent on the Zariski site, vanishes on affines and computes the same cohomology}
\author{The Stacks Project Authors}
\date{stacks-project:adequate.tex:local-snapshot-20261012; GNU FDL 1.2 or later}
\begin{document}
\maketitle
\noindent\textit{Editorial note.} All delivered mathematical text is the author's own, copied from the pinned Stacks Project chapter file adequate.tex; no statement or proof step was rewritten, reordered or removed.\par\medskip
\noindent\textit{Editorial note.} The counted spans carry navigation stubs for the references leaving them; no bibliography citation occurs in the spans.\par\medskip
\noindent\textit{Editorial note.} The preamble span points at the chapter's own preamble.tex shipped with the frozen source, used verbatim as staged.\par\medskip
\noindent\textit{Allowed standard prerequisites (editorial).}\begin{itemize}
\item The flat base change property of adequate modules and their locality. Same chapter, Lemmas lemma-adequate-flat and lemma-adequate-local, invoked at proof lines 1914-1921. Give F(U) = F(V) tensor O(U) for flat U -> V, hence quasi-coherence of the Zariski restriction.
\item Cohomology of a quasi-coherent module on an affine scheme vanishes in positive degrees, and the comparison between Zariski and site cohomology. The standard Cech computation for quasi-coherent modules, used in the closing paragraph of the proof. Gives the vanishing and the comparison statement.
\end{itemize}\medskip
\ifdefined\hypertarget\hypertarget{sourcepointer-1}{}\fi
\noindent\textit{Source reference.} Original lemma 03FI: sources/stacks-chow/tags.txt, source lines 4493--4493. Pointer outside the counted span, invoked once for context; the load-bearing argument of the counted statement is the author's proof reproduced in full in the proof span.\par\medskip
\ifdefined\hypertarget\hypertarget{sourcepointer-2}{}\fi
\noindent\textit{Source reference.} Original proposition 03DW: sources/stacks-chow/tags.txt, source lines 4435--4435. Pointer outside the counted span, invoked once for context; the load-bearing argument of the counted statement is the author's proof reproduced in full in the proof span.\par\medskip
\ifdefined\hypertarget\hypertarget{sourcepointer-3}{}\fi
\noindent\textit{Source reference.} Original lemma 06UW: sources/stacks-adequate/pinned-original-source.tex, source lines 442--442. Pointer outside the counted span, invoked once for context; the load-bearing argument of the counted statement is the author's proof reproduced in full in the proof span.\par\medskip
\ifdefined\hypertarget\hypertarget{sourcepointer-4}{}\fi
\noindent\textit{Source reference.} Original lemma 06VL: sources/stacks-adequate/pinned-original-source.tex, source lines 1841--1841. Pointer outside the counted span, invoked once for context; the load-bearing argument of the counted statement is the author's proof reproduced in full in the proof span.\par\medskip
\ifdefined\hypertarget\hypertarget{sourcepointer-5}{}\fi
\noindent\textit{Source reference.} Original lemma 06VI: sources/stacks-adequate/pinned-original-source.tex, source lines 1707--1707. Pointer outside the counted span, invoked once for context; the load-bearing argument of the counted statement is the author's proof reproduced in full in the proof span.\par\medskip
\ifdefined\hypertarget\hypertarget{sourcepointer-6}{}\fi
\noindent\textit{Source reference.} Original lemma 03F9: sources/stacks-chow/tags.txt, source lines 4484--4484. Pointer outside the counted span, invoked once for context; the load-bearing argument of the counted statement is the author's proof reproduced in full in the proof span.\par\medskip
\begin{lemma}
\label{lemma-same-cohomology-adequate}
Let $S$ be a scheme. Let $\mathcal{F}$ be an adequate
$\mathcal{O}$-module on $(\Sch/S)_\tau$.
\begin{enumerate}
\item The restriction $\mathcal{F}|_{S_{Zar}}$ is a quasi-coherent
$\mathcal{O}_S$-module on the scheme $S$.
\item The restriction $\mathcal{F}|_{S_\etale}$ is the
quasi-coherent module associated to $\mathcal{F}|_{S_{Zar}}$.
\item For any affine scheme $U$ over $S$ we have $H^q(U, \mathcal{F}) = 0$
for all $q > 0$.
\item There is a canonical isomorphism
$$
H^q(S, \mathcal{F}|_{S_{Zar}}) =
H^q((\Sch/S)_\tau, \mathcal{F}).
$$
\end{enumerate}
\end{lemma}

\begin{proof}
By
Lemma \ref{lemma-adequate-flat}
and
Lemma \ref{lemma-adequate-local}
we see that for any flat morphism of affines $U \to V$ over $S$
we have
$\mathcal{F}(U) = \mathcal{F}(V) \otimes_{\mathcal{O}(V)} \mathcal{O}(U)$.
This works in particular if $U \subset V \subset S$ are affine opens of
$S$, hence $\mathcal{F}|_{S_{Zar}}$ is quasi-coherent.
Thus (1) holds.

\medskip\noindent
Let $S' \to S$ be an \'etale morphism of schemes.
Then for $U \subset S'$ affine open mapping into an affine open
$V \subset S$ we see that
$\mathcal{F}(U) = \mathcal{F}(V) \otimes_{\mathcal{O}(V)} \mathcal{O}(U)$
because $U \to V$ is \'etale, hence flat. Therefore
$\mathcal{F}|_{S'_{Zar}}$ is the pullback of $\mathcal{F}|_{S_{Zar}}$.
This proves (2).

\medskip\noindent
We are going to apply
Cohomology on Sites,
Lemma \ref{sites-cohomology-lemma-cech-vanish-collection}
to the site $(\Sch/S)_\tau$ with
$\mathcal{B}$ the set of affine schemes over $S$ and
$\text{Cov}$ the set of standard affine $\tau$-coverings.
Assumption (3) of the lemma is satisfied by
Descent, Lemma \ref{descent-lemma-standard-covering-Cech}
and
Lemma \ref{lemma-adequate-fpqc}
for the case of a covering by a single affine.
Hence we conclude that $H^p(U, \mathcal{F}) = 0$ for every
affine scheme $U$ over $S$. This proves (3).
In exactly the same way as in the proof of
Descent, Proposition \ref{descent-proposition-same-cohomology-quasi-coherent}
this implies the equality of cohomologies (4).
\end{proof}

\end{document}
